% --- CORRECCIÓN IMPORTANTE EN LA PRIMERA LÍNEA ---
% Agregamos "xcolor=table" para que funcione \rowcolor sin errores
\documentclass[aspectratio=169, 10pt, xcolor=table]{beamer}

% --- Configuración del Tema (estilo profesional como el ejemplo) ---
\usetheme{Madrid}
\usecolortheme{dolphin}

% --- Paquetes necesarios ---
\usepackage[utf8]{inputenc}
\usepackage[spanish]{babel}
\usepackage{amsmath, amssymb}
\usepackage{booktabs}
\usepackage{graphicx}
\usepackage{listings}
\usepackage{tikz}
\usetikzlibrary{arrows.meta, positioning, shapes.geometric}

% --- Ruta de Imágenes (crea la carpeta "imagenes" y coloca ahí tus archivos) ---
\graphicspath{{./imagenes/}}

% --- Estilo para código Python ---
\definecolor{codegreen}{rgb}{0,0.6,0}
\definecolor{codegray}{rgb}{0.5,0.5,0.5}
\definecolor{codepurple}{rgb}{0.58,0,0.82}
\definecolor{backcolour}{rgb}{0.95,0.95,0.92}
\lstset{
    language=Python,
    backgroundcolor=\color{backcolour},
    commentstyle=\color{codegreen},
    keywordstyle=\color{blue},
    numberstyle=\tiny\color{codegray},
    stringstyle=\color{codepurple},
    basicstyle=\ttfamily\scriptsize,
    breaklines=true,
    frame=single,
    showstringspaces=false
}

% --- Pie de página personalizado ---
\makeatletter

\setbeamertemplate{footline}{
  \leavevmode%
  \hbox{%
  \begin{beamercolorbox}[wd=.333333\paperwidth,ht=2.25ex,dp=1ex,center]{author in head/foot}%
    \usebeamerfont{author in head/foot}\insertshortauthor
  \end{beamercolorbox}%
  \begin{beamercolorbox}[wd=.333333\paperwidth,ht=2.25ex,dp=1ex,center]{title in head/foot}%
    \usebeamerfont{title in head/foot}Sesión 13
  \end{beamercolorbox}%
  \begin{beamercolorbox}[wd=.333333\paperwidth,ht=2.25ex,dp=1ex,right]{date in head/foot}%
    \usebeamerfont{date in head/foot}Machine Learning \hfill \insertframenumber/\inserttotalframenumber \hspace*{0.3cm}
  \end{beamercolorbox}}%
  \vskip0pt%
}

\makeatother

% --- Datos de la portada ---
\title[SESIÓN 13]{\textbf{Sesión 13}}
\subtitle{Interpretabilidad de Modelos: \\ Explicando las decisiones de la caja negra}
\author[Prof. C. Sánchez]{Carlos César Sánchez Coronel}
\institute{\textbf{MACHINE LEARNING}}
\date{}

% Logo opcional (descomentar si tienes la imagen)
% \titlegraphic{\includegraphics[height=1.5cm]{logo.png}}

% --- Eliminar la barra de navegación (opcional) ---
\setbeamertemplate{navigation symbols}{}

% --- Índice al inicio de cada sección ---
\AtBeginSection[]{
  \begin{frame}
    \frametitle{Contenido}
    \tableofcontents[currentsection]
  \end{frame}
}

% ============================================================================
% INICIO DEL DOCUMENTO
% ============================================================================
\begin{document}

% --- Portada ---
\begin{frame}
    \titlepage
\end{frame}

% --- Índice general ---
\begin{frame}{Contenido}
    \tableofcontents
\end{frame}

% ============================================================================
% SECCIÓN 1: INTRODUCCIÓN
% ============================================================================
\section{Introducción}

\begin{frame}{La necesidad de explicaciones}
    \begin{itemize}
        \item Modelos complejos (XGBoost, redes neuronales) son ``cajas negras''.
        \item En ámbitos críticos (salud, finanzas, justicia) la transparencia es obligatoria.
        \item Regulaciones como el GDPR otorgan el ``derecho a una explicación''.
        \item La interpretabilidad también ayuda a:
        \begin{itemize}
            \item Depurar modelos (detectar sesgos, errores).
            \item Generar confianza en usuarios y stakeholders.
            \item Descubrir nuevo conocimiento.
        \end{itemize}
    \end{itemize}
    % Basado en Pregunta 1
    % IMAGEN: Esquema caja negra vs transparente
\end{frame}

\begin{frame}{Tipos de interpretabilidad}
    \begin{columns}[t]
        \column{0.5\textwidth}
        \textbf{Interpretabilidad global}
        \begin{itemize}
            \item Comportamiento promedio del modelo.
            \item ¿Qué variables son más importantes en general?
            \item ¿Cómo afectan en promedio?
        \end{itemize}
        \column{0.5\textwidth}
        \textbf{Interpretabilidad local}
        \begin{itemize}
            \item Explicación de una predicción individual.
            \item ¿Por qué este cliente fue rechazado?
            \item Necesario para dar explicaciones personalizadas.
        \end{itemize}
    \end{columns}
    \vspace{0.3cm}
    \begin{block}{Métodos}
        \begin{itemize}
            \item \textbf{Intrínsecos}: el modelo es interpretable por sí mismo (regresión lineal, árboles pequeños).
            \item \textbf{Post-hoc}: se aplican después de entrenar (SHAP, LIME, PDP).
        \end{itemize}
    \end{block}
    % Basado en Pregunta 1
\end{frame}

% ============================================================================
% SECCIÓN 2: IMPORTANCIA DE CARACTERÍSTICAS EN ÁRBOLES
% ============================================================================
\section{Importancia de Características en Árboles}

\begin{frame}{Importancia por impureza (MDI)}
    \begin{block}{Definición}
        En árboles y Random Forest, la importancia de una característica se calcula sumando la reducción de impureza (Gini o entropía) en todos los nodos donde aparece, ponderada por la proporción de muestras.
        \[
        \text{Imp}(j) = \frac{\sum_{t \in T} \Delta i(t) \cdot n(t)}{\sum_{t \in T} n(t)}
        \]
    \end{block}
    \begin{itemize}
        \item \textbf{Ventaja}: rápida, incluida en la mayoría de librerías.
        \item \textbf{Limitaciones}: sesgo hacia variables numéricas o con muchas categorías; no indica dirección del efecto; puede ser engañosa con correlaciones.
    \end{itemize}
    % Basado en Pregunta 3
\end{frame}

\begin{frame}{Importancia por permutación}
    \begin{block}{Idea}
        Si una variable es importante, permutar sus valores rompe la relación con la variable objetivo, aumentando el error.
        \[
        \text{Imp}(j) = \frac{1}{K}\sum_{k=1}^{K} (\text{error}_{\text{perm}} - \text{error}_{\text{original}})
        \]
    \end{block}
    \begin{itemize}
        \item Más fiable que MDI, modelo-agnóstico, independiente de la escala.
        \item Método recomendado en scikit-learn.
    \end{itemize}
    % Basado en Pregunta 2
\end{frame}

\begin{frame}[fragile]{Importancia por permutación en Python}
    \begin{lstlisting}
from sklearn.inspection import permutation_importance
from sklearn.ensemble import RandomForestClassifier

model = RandomForestClassifier().fit(X_train, y_train)
result = permutation_importance(model, X_test, y_test, n_repeats=10)
sorted_idx = result.importances_mean.argsort()
    \end{lstlisting}
\end{frame}

% ============================================================================
% SECCIÓN 3: PARTIAL DEPENDENCE PLOTS (PDP) E ICE
% ============================================================================
\section{PDP e ICE}

\begin{frame}{Partial Dependence Plot (PDP)}
    \begin{block}{Definición}
        Muestra cómo cambia la predicción promedio al variar una (o dos) características, marginalizando sobre el resto.
        \[
        \widehat{f}_{x_s}(x_s) = \frac{1}{n}\sum_{i=1}^{n} f(x_s, x_c^{(i)})
        \]
    \end{block}
    \begin{itemize}
        \item \textbf{Interpretación}: relación monótona, lineal, o compleja.
        \item \textbf{Limitación}: asume independencia entre $x_s$ y $x_c$; puede promediar regiones irreales si hay correlación.
    \end{itemize}
    % Basado en Pregunta 4
    % IMAGEN: PDP de ejemplo
\end{frame}

\begin{frame}{Individual Conditional Expectation (ICE)}
    \begin{block}{Idea}
        Para cada instancia, se mantienen fijas sus otras variables y se varía $x_s$, obteniendo una curva por observación.
    \end{block}
    \begin{itemize}
        \item El PDP es el promedio de todas las curvas ICE.
        \item Si las curvas ICE no son paralelas, hay interacciones.
        \item Útil para detectar heterogeneidad en el efecto de una variable.
    \end{itemize}
    % Basado en Pregunta 5
\end{frame}

\begin{frame}[fragile]{PDP e ICE con scikit-learn}
    \begin{lstlisting}
from sklearn.inspection import PartialDependenceDisplay

# PDP
PartialDependenceDisplay.from_estimator(model, X_train,
    features=['duration', 'amount'], kind='average')

# ICE
PartialDependenceDisplay.from_estimator(model, X_train,
    features=['duration'], kind='individual')
    \end{lstlisting}
\end{frame}

% ============================================================================
% SECCIÓN 4: SHAP (SHapley Additive exPlanations)
% ============================================================================
\section{SHAP}

\begin{frame}{Valores de Shapley en teoría de juegos}
    \begin{block}{Fórmula general}
        \[
        \phi_j = \sum_{S \subseteq P \setminus \{j\}} \frac{|S|!(p - |S| - 1)!}{p!} [v(S \cup \{j\}) - v(S)]
        \]
        donde $v(S)$ es la predicción usando solo las características en $S$.
    \end{block}
    \begin{itemize}
        \item Asigna a cada característica su contribución justa a la predicción.
        \item Propiedades deseables: eficiencia, simetría, linealidad, dummy.
    \end{itemize}
    % Basado en Pregunta 6
\end{frame}

\begin{frame}{SHAP en la práctica}
    \begin{itemize}
        \item \textbf{Kernel SHAP}: modelo-agnóstico, combina LIME con propiedades Shapley (lento).
        \item \textbf{Tree SHAP}: específico para árboles, rápido y exacto (aprovecha la estructura de los árboles).
        \item \textbf{Deep SHAP}: para redes neuronales.
    \end{itemize}
    \vspace{0.3cm}
    \begin{block}{Interpretación}
        \begin{itemize}
            \item Valor base (expected value): predicción media en el conjunto de entrenamiento.
            \item SHAP values: contribuciones de cada característica.
            \item Suma: $\text{valor base} + \sum \phi_j = \text{predicción}$ (propiedad de eficiencia).
        \end{itemize}
    \end{block}
    % Basado en Pregunta 7
\end{frame}

\begin{frame}{Visualizaciones SHAP}
    \begin{itemize}
        \item \textbf{Summary plot}: importancia global + dirección del efecto.
        \item \textbf{Dependence plot}: relación entre el valor de la variable y su impacto SHAP; puede colorearse por otra variable para detectar interacciones.
        \item \textbf{Force plot / Waterfall plot}: explicación local para una instancia.
    \end{itemize}
    % Basado en Preguntas 8 y 9
    % IMAGEN: Ejemplos de summary plot, force plot
\end{frame}

\begin{frame}[fragile]{TreeExplainer en Python}
    \begin{lstlisting}
import shap

explainer = shap.TreeExplainer(model)
shap_values = explainer.shap_values(X_test)

# Summary plot
shap.summary_plot(shap_values, X_test, feature_names=feature_names)

# Force plot para una instancia
shap.force_plot(explainer.expected_value, shap_values[0,:],
                X_test.iloc[0,:])
    \end{lstlisting}
\end{frame}

% ============================================================================
% SECCIÓN 5: LIME (Local Interpretable Model-agnostic Explanations)
% ============================================================================
\section{LIME}

\begin{frame}{LIME: idea fundamental}
    \begin{block}{Pasos}
        \begin{enumerate}
            \item Generar muestras aleatorias alrededor de la instancia a explicar.
            \item Obtener predicciones del modelo original para esas muestras.
            \item Ponderar las muestras según su proximidad (kernel exponencial).
            \item Ajustar un modelo interpretable (lineal) en esa vecindad.
            \item Los coeficientes del modelo simple son la explicación.
        \end{enumerate}
    \end{block}
    % Basado en Pregunta 9
\end{frame}

\begin{frame}{Ventajas y limitaciones de LIME}
    \begin{itemize}
        \item \textbf{Ventajas}: agnóstico al modelo, rápido, intuitivo.
        \item \textbf{Limitaciones}: inestable (depende del muestreo aleatorio), difícil definir la vecindad.
    \end{itemize}
    % Basado en Pregunta 10
\end{frame}

\begin{frame}[fragile]{LIME en Python}
    \begin{lstlisting}
from lime.lime_tabular import LimeTabularExplainer

explainer = LimeTabularExplainer(X_train.values,
                                 feature_names=feature_names,
                                 class_names=['bueno', 'malo'],
                                 mode='classification')
exp = explainer.explain_instance(X_test.iloc[0].values,
                                 model.predict_proba)
exp.show_in_notebook()
    \end{lstlisting}
\end{frame}

% ============================================================================
% SECCIÓN 6: COMPARACIÓN DE MÉTODOS
% ============================================================================
\section{Comparación de Métodos}

\begin{frame}{Comparación: SHAP vs LIME vs PDP vs Importancia}
    \begin{table}
        \centering
        \begin{tabular}{l|c|c|c}
            \toprule
            \textbf{Método} & \textbf{Ámbito} & \textbf{Modelo-agnóstico} & \textbf{Velocidad} \\
            \midrule
            Importancia (permutación) & Global & Sí & Rápida \\
            PDP & Global & Sí & Media \\
            ICE & Local/Global & Sí & Media \\
            SHAP (Tree) & Local/Global & No & Rápida \\
            SHAP (Kernel) & Local/Global & Sí & Lenta \\
            LIME & Local & Sí & Media \\
            \bottomrule
        \end{tabular}
    \end{table}
    % Basado en teoría y Pregunta 11
\end{frame}

\begin{frame}{¿Cuándo usar cada uno?}
    \begin{itemize}
        \item \textbf{Importancia por permutación}: ranking global rápido y fiable.
        \item \textbf{PDP/ICE}: entender la relación funcional de una variable.
        \item \textbf{SHAP}: explicaciones locales consistentes + resumen global (estándar de facto).
        \item \textbf{LIME}: alternativa rápida cuando no se puede usar SHAP, o para modelos no árbol.
    \end{itemize}
    % Basado en Pregunta 11
\end{frame}

% ============================================================================
% SECCIÓN 7: APLICACIONES EN LA INDUSTRIA
% ============================================================================
\section{Aplicaciones en la Industria}

\begin{frame}{Finanzas: scoring crediticio}
    \begin{itemize}
        \item Cliente rechazado exige explicación (GDPR).
        \item Usar SHAP force plot para mostrar contribuciones.
        \item Generar informe en lenguaje natural con los factores principales.
    \end{itemize}
    % Basado en Pregunta 1
\end{frame}

\begin{frame}{Salud: diagnóstico asistido}
    \begin{itemize}
        \item Médico necesita entender por qué el modelo predice alto riesgo de diabetes.
        \item PDP de glucosa, IMC, edad muestra efectos.
        \item SHAP local para un paciente específico.
    \end{itemize}
    % Basado en Pregunta 14
\end{frame}

\begin{frame}{Recursos humanos: detección de sesgos}
    \begin{itemize}
        \item Analizar valores SHAP por grupo demográfico.
        \item Calcular métricas de equidad (disparate impact, igualdad de oportunidades).
        \item Identificar variables proxy.
    \end{itemize}
    % Basado en Pregunta 15
\end{frame}

% ============================================================================
% SECCIÓN 8: CASO INTEGRADO - GERMAN CREDIT
% ============================================================================
\section{Caso Integrado: German Credit}

\begin{frame}{Datos y modelo}
    \begin{itemize}
        \item Dataset German Credit (UCI): 1000 solicitudes, 20 variables, objetivo binario (bueno/malo).
        \item Modelo: XGBoost con AUC 0.78.
        \item Variables clave: duration, amount, age, purpose, other_installment_plans.
    \end{itemize}
    % Basado en Pregunta 19
\end{frame}

\begin{frame}{Interpretabilidad global}
    \begin{itemize}
        \item Importancia por permutación: duration, amount, age.
        \item PDP de duration: la probabilidad de malo aumenta con la duración.
        \item SHAP summary plot: valores altos de duration y amount contribuyen al riesgo (SHAP positivo); age joven también.
    \end{itemize}
    % Basado en Pregunta 19
\end{frame}

\begin{frame}{Explicación local (cliente rechazado)}
    \begin{block}{Datos del cliente}
        \begin{itemize}
            \item duration = 48 meses
            \item amount = 5000 €
            \item age = 25 años
            \item purpose = coche
            \item other_installment_plans = yes
        \end{itemize}
        Predicción: 0.82 (riesgo alto).
    \end{block}
    \begin{block}{SHAP force plot (simulado)}
        \begin{itemize}
            \item Valor base: 0.35
            \item duration=48: +0.25
            \item amount=5000: +0.15
            \item age=25: +0.10
            \item other_installment_plans=yes: +0.07
            \item resto: -0.05
            \item Suma: 0.35+0.25+0.15+0.10+0.07-0.05 = 0.87 (aprox.)
        \end{itemize}
    \end{block}
\end{frame}

\begin{frame}{Comunicación al cliente}
    \begin{quote}
        "Su solicitud fue rechazada principalmente porque la duración del crédito (48 meses) es superior a la media, lo que incrementa el riesgo. También influyen el monto solicitado (5000 €), su edad (25 años) y el hecho de tener otros planes de pago. Estos factores combinados resultan en una alta probabilidad de incumplimiento."
    \end{quote}
    % Basado en Pregunta 19
\end{frame}

% ============================================================================
% SECCIÓN 9: LIMITACIONES Y BUENAS PRÁCTICAS
% ============================================================================
\section{Limitaciones y Buenas Prácticas}

\begin{frame}{Limitaciones de los métodos}
    \begin{itemize}
        \item \textbf{SHAP}: con correlaciones extremas, las contribuciones pueden repartirse de manera inestable; no implica causalidad.
        \item \textbf{LIME}: inestable por muestreo; ejecutar varias veces para verificar consistencia.
        \item \textbf{PDP}: puede promediar regiones irreales si hay correlaciones (usar ALE plots).
        \item \textbf{Importancia}: no indica dirección del efecto.
    \end{itemize}
    % Basado en teoría y Pregunta 13
\end{frame}

\begin{frame}{Buenas prácticas}
    \begin{itemize}
        \item Validar explicaciones con expertos del dominio.
        \item No confundir correlación con causalidad.
        \item Usar múltiples métodos para obtener una visión completa.
        \item Para cumplimiento regulatorio, documentar el proceso de interpretabilidad.
    \end{itemize}
\end{frame}

% ============================================================================
% SECCIÓN 10: CONEXIÓN Y RESUMEN
% ============================================================================
\section{Conexión y Resumen}

\begin{frame}{Conexión con el resto del curso}
    \begin{itemize}
        \item Importancia de características se basa en métricas de sesiones anteriores.
        \item SHAP utiliza teoría de juegos (matemáticas discretas).
        \item LIME usa regresión lineal local (modelos lineales).
        \item Interpretabilidad es clave para MLOps y gobernanza de modelos (sesión 14).
    \end{itemize}
\end{frame}

\begin{frame}{Lo que hemos aprendido}
    \begin{exampleblock}{Conceptos clave}
        \begin{itemize}
            \item Interpretabilidad global vs local.
            \item Importancia por permutación (robusta) vs MDI (sesgada).
            \item PDP e ICE para relaciones parciales.
            \item SHAP: valores de Shapley, TreeExplainer, visualizaciones.
            \item LIME: explicaciones locales rápidas.
            \item Aplicaciones en finanzas, salud, detección de sesgos.
            \item Caso práctico: German Credit.
        \end{itemize}
    \end{exampleblock}
\end{frame}

% --- Diapositiva final ---
\begin{frame}
    \centering
    \Huge \textbf{¡Gracias!}

\end{frame}

\end{document}